\section{Introduction}\label{sec:introduction}

Pure infiniteness expresses the ability of positive elements of a
\(C^*\)-algebra to contain several copies of themselves in Cuntz
comparison.  For a nonsimple algebra, an individual comparison and
the simultaneous comparison of two elements are different demands:
the latter must also control a mixed entry.  We prove that proper
infiniteness of every positive element already provides this
compatibility.  We then give two ways to remove a fixed finite
amplification, one for exact algebras and an independent construction
for nuclear algebras.  In the separable nuclear case the result yields
absorption of the Cuntz algebra \(\Oinf\).

\subsection{Infiniteness notions and main results}

All \(C^*\)-algebras are complex and may be nonunital.  Let \(\Kc\)
denote the compact operators on an infinite-dimensional separable
Hilbert space.  For \(u,v\in(C\otimes\Kc)_+\), write \(u\precsim v\)
if some sequence \(x_j\in C\otimes\Kc\) satisfies
\(x_j^*vx_j\to u\) in norm, and write \(u\CuEq v\) if comparison holds
in both directions.  This is stable Cuntz comparison
\cite{Cuntz1978}.  Finite matrices over \(C\) occupy the standard
corners of \(C\otimes\Kc\); comparisons between different sizes use
zero padding in the other coordinates.  Write \(h^{\oplus n}\) for
the direct sum of \(n\) copies of \(h\).  A positive element \(h\) is
\emph{properly infinite} if \(h\oplus h\precsim h\).  We include zero
as properly infinite.  For \(h\in C_+\), let
\[
 \Her(h)=\overline{hCh},\qquad
 \Id(h)=\text{the closed two-sided ideal of \(C\) generated by \(h\)}.
\]

We will use the following two comparison conditions:
\begin{align}
 &h\text{ is properly infinite for every }h\in C_+,
 \tag{P}\label{eq:property-P}\\
 &h^{\oplus n}\text{ is properly infinite for every }h\in C_+,
 \quad\text{for one fixed integer }n\ge1.
 \tag{\(F_n\)}\label{eq:property-F}
\end{align}
In \(\propF\), the same \(n\) is chosen before all positive elements
are tested.  It therefore applies to every scalar spectral cutoff
used later.  The premise concerns scalar positives in \(C\), with
their indicated amplifications.

Condition \(\propP\) is equivalent to pure infiniteness in the
nonsimple sense of Kirchberg--R{\o}rdam
\cite[Theorem~4.16]{KR2000}; the equivalence is also recalled in
\cite[Section~21]{STW99problems}.  The existence of an
\(n\) satisfying \(\propF\) is \emph{weak pure infiniteness} in the
fixed-amplification sense.  We use the following formulation of
\emph{strong pure infiniteness}, valid for arbitrary \(C^*\)-algebras
\cite[Definition~3.1]{KS}: for every \(x,y\in C_+\) and
\(\varepsilon>0\), there are \(s,t\in C\) such that
\begin{equation}\label{eq:SPI-definition}
 \|s^*x^2s-x^2\|<\varepsilon,\qquad
 \|t^*y^2t-y^2\|<\varepsilon,\qquad
 \|s^*xy t\|<\varepsilon.
\end{equation}
The coefficients in this formulation have no prescribed norm bound.

Kirchberg--R{\o}rdam introduced the weak and strong notions in
\cite{KR2002}.  Strong pure infiniteness implies ordinary pure
infiniteness, and \(\propP\) implies \(\propF\) with \(n=1\).
The reverse implications form the comparison question in
Kirchberg--R{\o}rdam's Question~9.5, as recalled in
\cite[Section~21, Problem~LXXII]{STW99problems}.  In version~2 of that problem collection, the question asks whether
the three notions agree in general and in the nuclear setting.
Our first theorem proves the implication from ordinary to strong
pure infiniteness for arbitrary complex \(C^*\)-algebras, in the
following more flexible form.

\begin{theorem}\label{thm:individual-to-strong}
Let \(C\) be a complex \(C^*\)-algebra satisfying \(\propP\).
For all \(a,b\in C_+\), \(c\in C\), and \(\varepsilon>0\), there are
\(s,t\in C\) such that
\begin{equation}\label{eq:main-diagonalization}
 \|s^*as-a\|<\varepsilon,\qquad
 \|t^*bt-b\|<\varepsilon,\qquad
 \|s^*ct\|<\varepsilon.
\end{equation}
Consequently \(C\) is strongly purely infinite.
\end{theorem}

The cross term \(c\) is arbitrary and need not be a product of the
diagonal inputs.  Taking \((a,b,c)=(x^2,y^2,xy)\) gives the stated
consequence.  Arbitrary-entry diagonalization is a known equivalent
form of strong pure infiniteness
\cite[Remarks~3.2(i)--(ii)]{KS}; our proof obtains it directly from
\(\propP\).  It also gives coefficient bounds for positive
contraction inputs that are independent of the requested mixed-entry
accuracy; see Proposition~\ref{prop:bounded-diagonalization}.

Recall that a \(C^*\)-algebra is exact if its minimal tensor product
preserves short exact sequences.  Exactness supplies the additional
uniformity needed to remove the fixed amplification.

\begin{theorem}\label{thm:exact-infiniteness}
Let \(C\) be an exact complex \(C^*\)-algebra.  Suppose that there is
one integer \(n\ge1\) such that \(h^{\oplus n}\) is properly infinite
for every \(h\in C_+\).  Then every positive element of \(C\) is
properly infinite.  Consequently \(C\) is strongly purely infinite.
\end{theorem}

These theorems give a positive answer to the ordinary-to-strong part
of Problem~LXXII in general and to the full comparison question for
exact algebras, hence also for nuclear algebras.  The weak-to-ordinary implication
for arbitrary nonexact algebras remains outside these conclusions.
The statements include nonseparable and nonunital algebras.  Section~\ref{sec:nuclear}
gives an independent proof of the first conclusion of
Theorem~\ref{thm:exact-infiniteness} under nuclearity.  Its statement
is a restriction of that theorem, while its finite-tuple obstruction
and merger of hereditary witnesses provide a different method.
The same method extends to exact algebras by the deeper theorem that
exactness passes to quotients; see Remark~\ref{rem:tuple-exact}.
The direct exact proof avoids that permanence theorem.

Let \(\Oinf\) be the Cuntz algebra generated by countably many
isometries with mutually orthogonal ranges.  Kirchberg--R{\o}rdam
proved the equivalence between strong pure infiniteness and
\(\Oinf\)-absorption for separable nuclear algebras that are stable
or have an approximate unit of projections \cite{KR2002}.
Kirchberg's later central-sequence results cover general nonunital
separable nuclear algebras
\cite[Proposition~4.4 and Remark~4.6]{Kirchberg2006}; see also the
paragraph following \cite[Definition~3.1]{KS}.  At that scope they
give the following tensorial consequence.

\begin{corollary}\label{cor:absorption}
Let \(A\) be a separable nuclear complex \(C^*\)-algebra.  If there is
one integer \(n\ge1\) such that \(h^{\oplus n}\) is properly infinite
for every \(h\in A_+\), then
\[
 A\cong A\otimes_{\min}\Oinf.
\]
\end{corollary}

For separable nuclear algebras, therefore, weak pure infiniteness,
pure infiniteness, strong pure infiniteness, and \(\Oinf\)-absorption
are equivalent.  The reverse implication from absorption to strong
pure infiniteness is due to Kirchberg--R{\o}rdam
\cite{KR2002}; see also the discussion after
\cite[Definition~3.1]{KS}.  Thus a single finite amplification test
on positive elements characterizes this tensorial regularity even
when the algebra is nonsimple or nonunital.

For example, the cone \(C_0((0,1])\otimes\Oinf\) is separable
nuclear and absorbs \(\Oinf\) by the self-absorption of \(\Oinf\).
It is nonsimple, nonunital, and has no nonzero projections: the norm
of a continuous projection-valued function takes only the values
zero and one and must tend to zero at the missing endpoint.

\subsection{Earlier results and the role of ideal structure}

The comparison question has been settled under several structural
conditions.  Kirchberg--R{\o}rdam treat simple, real-rank-zero, and
approximately divisible cases \cite{KR2002}; the real-rank-zero
conclusion is also recalled in \cite[Section~21]{STW99problems}.
Blanchard--Kirchberg prove that ordinary pure infiniteness implies
strong pure infiniteness when the primitive ideal space is Hausdorff
\cite[Theorem~5.8]{BlanchardKirchberg2004}.  R{\o}rdam's proof of
\cite[Theorem~5.2]{Rordam2004Z} gives the same implication for
\(\mathcal Z\)-absorbing purely infinite algebras, where
\(\mathcal Z\) is the Jiang--Su algebra.  These two implications do
not require exactness or nuclearity.

Pasnicu--R{\o}rdam establish equivalence of the three infiniteness
notions for algebras with the ideal property, meaning that projections
separate ideals \cite[Proposition~2.14]{PasnicuRordam2007}.
For algebras of topological dimension zero, Pasnicu--Rouzbehani prove
the implication from ordinary to strong pure infiniteness
\cite[Theorems~3.10(i) and~3.13]{PasnicuRouzbehani2025}.
Here topological dimension zero means that the primitive ideal space
has a basis of compact open sets.  Theorem~\ref{thm:individual-to-strong}
requires no restriction on this space or on the ideal structure.

The weak-to-ordinary implication is closely related to the problem of
finding sufficiently many orthogonal equivalent elements in a
hereditary subalgebra.  Kirchberg--R{\o}rdam showed that the Global
Glimm Halving property permits this implication
\cite[Proposition~4.15]{KR2002}; see
\cite[Theorem~39]{STW99problems}.  An equivalent formulation of the
Global Glimm Property is that, for every \(a\in C_+\) and
\(\varepsilon>0\), there is \(x\in\Her(a)\) with \(x^2=0\) such
that \((a-\varepsilon)_+\) belongs to the closed two-sided ideal
generated by \(x\) \cite[Remark~3.2 and Theorem~3.6(3)]{TV2023}.
Elliott--Rouzbehani prove weak-to-ordinary pure infiniteness in
topological dimension zero \cite{ER2023}.  Ng--Thiel--Vilalta recover
that implication through the Global Glimm Property in a proof that
includes nonseparable algebras
\cite[Theorem~2.3 and Corollary~2.5]{NTV2026}.
The initial and range squares of a square-zero element are
orthogonal and Cuntz equivalent.  Our exact argument constructs such
elements while retaining a nonzero image throughout a prescribed
set of primitive quotients.  The relevant sets are the norm slices
\(\{P\in\Prim(C):\|h+P\|\ge r\}\), where
\(h\in C_+\) and \(r>0\); they are quasi-compact even when
\(\Prim(C)\) is not Hausdorff.  Repeated halving on the same slice
then supplies enough copies to remove the fixed amplification.  The nuclear argument instead combines
representation-dependent positive elements into one element that
covers a prescribed cutoff.

\subsection{Proof strategy and the independent routes}

Two tasks drive the argument.  To pass from \(\propP\) to strong
pure infiniteness, one must choose two hereditary corners that retain
the relevant parts of their inputs while making the mixed term small.
To remove \(n\) from \(\propF\), one must first produce sufficiently
many equivalent orthogonal pieces inside a scalar hereditary
subalgebra.  Figure~\ref{fig:proof-routes} separates these tasks and
the two methods for the second one.

\begin{figure}[htbp]
\centering
\small
\fbox{\strut Comparison and cutoff transport (Section~\ref{sec:foundations})}
\par\medskip
\newcommand{\routebox}[1]{\fbox{\parbox[c][13mm][c]{0.27\textwidth}{\centering #1}}}
\begin{tabular}{@{}c@{\hspace{3mm}}c@{\hspace{3mm}}c@{}}
\textbf{Arbitrary algebra, \(\propP\)}&
\textbf{Nuclear algebra, \(\propF\)}&
\textbf{Exact algebra, \(\propF\)}\\[3pt]
\routebox{Graph and four entries\\Sections~\ref{sec:graph}--\ref{sec:four-variables}}&
\routebox{Graph and finite tuples\\Sections~\ref{sec:graph}, \ref{sec:nuclear}}&
\routebox{Uniform perturbation\\Section~\ref{sec:exactness}}\\
\(\Downarrow\)&\(\Downarrow\)&\(\Downarrow\)\\
\routebox{Products and iteration\\Section~\ref{sec:separation}}&
\routebox{Local elements and merger\\Section~\ref{sec:nuclear}}&
\routebox{Square-zero localization\\Section~\ref{sec:localization}}\\
\(\Downarrow\)&\(\Downarrow\)&\(\Downarrow\)\\
Strong pure infiniteness&\(\propP\)&\(\propP\)
\end{tabular}
\par\medskip
\(\propP\ \xRightarrow{\ \text{general route}\ }\
  \text{strong pure infiniteness}
  \ \xRightarrow{\ \text{separable nuclear}\ }\
  \Oinf\text{-absorption}\)
\caption{Proof routes and their hypotheses.  Comparison transport
supports all three routes.  Either output \(\propP\) can enter the
left-hand route.  Remark~\ref{rem:tuple-exact} extends the middle
route to exact algebras using quotient permanence.  Under exactness, Section~\ref{sec:quantitative}
gives an alternative passage from \(\propP\) to strong pure
infiniteness.  The final arrow is the external absorption criterion
used in Corollary~\ref{cor:absorption}.}
\label{fig:proof-routes}
\end{figure}


\paragraph{Comparison transport.}
Section~\ref{sec:foundations} supplies the common tool.  A Cuntz
comparison, after a spectral cutoff, can be implemented by a
rectangular partial isometry in the bidual.  Its products with the
cutoff hereditary algebra lie in the original algebra, and conjugation
preserves norms in every representation, also at matrix levels.
Thus one can place several orthogonal copies in a hereditary
subalgebra and transport an entire finite tuple through the same map.
These are stronger implementation data than a comparison of classes
alone.

\paragraph{From individual to simultaneous comparison.}
For \(h\ge0\), let \(e_h\) be its support projection in the universal
bidual.  Starting from positive contractions \(a,b\) and a contraction
\(c\), the first goal is to find positive contractions
\(p\in\Her(a)\), \(q\in\Her(b)\) with
\[
 \|e_pce_q\|<1,
\]
while requiring represented norm one of \(p\) or \(q\) whenever the
corresponding input has represented norm greater than \(1/2\).
This retention of represented norms is what later gives ideal
comparison.

Section~\ref{sec:graph} builds one corner as a graph inside several
orthogonal copies of a cutoff of \(a\).  The graph has finitely many
freely chosen algebra entries.  If the positive defect \(b(1-c^*p^2c)b\) vanishes in a
representation that detects the relevant part of \(b\), equality in
contraction inequalities forces all of
those entries to equal fixed measurements, computed before the entries
were chosen.  Under \(\propP\), Section~\ref{sec:four-variables}
excludes these equalities with four entries.  The first pair forces
an isometry with a nonzero defect, using the scaling-element
construction of Blackadar--Cuntz
\cite{BlackadarCuntz1982,PasnicuRordam2007}; the second pair is forced
both into and away from that defect.  Spectral cutoffs place the
relevant supports in actual operator ranges, and a directed family
of finite tests makes the contradiction work without separability.
Two ideal cutoffs then give the strict support estimate above.

Section~\ref{sec:separation} upgrades strict contraction to arbitrarily
small compression.  Uniformly bounded comparison rows show that
\(\propP\) passes to bounded products.  A product contradiction
therefore supplies one contraction factor below one, and iteration
preserves both represented norm signals while reducing the mixed
compression.  Finally, high spectral parts of \(a,b\) are compared
into the separated corners.  The coefficients have bounds determined
by the diagonal cutoffs, independent of the requested mixed-entry
accuracy.  This order of choices completes
Theorem~\ref{thm:individual-to-strong}.

\paragraph{A separate nuclear proof of amplification removal.}
Section~\ref{sec:nuclear} reuses the graph construction with a
different obstruction.  A finite tuple witnessing nonexactness has
linear matrix norms that no tuple in an exact algebra can realize
at every matrix level.  Finite-dimensional compressions nevertheless
approximate each fixed norm.  Under \(\propF\), one rectangular
carrier transports each entire compressed tuple into \(M_n(C)\).
If every choice failed to exclude the graph relations, their limiting
norms would occur in a quotient of \(M_n(C)\).  Nuclearity makes this
quotient nuclear, hence exact, giving a contradiction.  The matrix
norm and fixed-element quotient arguments are related to Pisier's
operator-space localization methods \cite[pp.~162--165]{Pisier1995}.

It remains to obtain the positive inputs that the graph needs and
combine them.  The fixed-amplification hypothesis and Kadison
transitivity \cite{Kadison1957} give, at each irreducible
representation detecting a positive \(u\), a positive contraction in \(\Her(u)\)
whose every spectral cutoff is properly infinite and whose norm in
that representation is one.  This element may
depend on the representation.  Repeated graph shrinking and a polar
deformation merge finitely many such elements while retaining all
their norm signals.  A finite ideal argument covers each cutoff of
\(u\) by one merged element and then compares two copies of that
cutoff back to \(u\).  This proves \(\propP\) under nuclearity,
independently of the exact argument.  The use of hereditary elements
to cover ideals is related to the filling-family criteria of
\cite[Theorem~1.1 and Definitions~4.2,~4.7]{KS}; here the merger supplies
the simultaneous norm control.

\paragraph{The exact proof and its quantitative completion.}
Section~\ref{sec:exactness} starts from the nonexact universal unital
\(C^*\)-algebra of one contraction.  Its finite matrix tests and its
embedded copy of \(C^*(\mathbb F_2)\) were studied by
Courtney--Sherman \cite{CS}, with Wassermann's nonexact free-group
example as a predecessor \cite{Wassermann1976}.  Finite-word
compression at all matrix levels turns nonexactness into one positive
error threshold: no contraction in a unital exact algebra can contain
approximately reducing copies of every finite matrix contraction.
The same copy must approximately intertwine both the contraction and
its adjoint.

Combining this obstruction with \(\propF\) gives a perturbation that
prevents cancellation simultaneously in every primitive quotient of a
prescribed norm slice.  Its lower norm bound is proportional to the
perturbation size, with a constant independent of the slice and of
\(n\).  Section~\ref{sec:localization} applies this estimate through
three finite coupling steps.  For a hereditary subalgebra nonzero
on a prescribed quasi-compact subset of \(\Prim(C)\), it constructs
a square-zero element nonzero at every point of that subset.  Repeated
halving on one fixed norm slice puts \(n\) copies of a positive
element below the input in Cuntz comparison; ideal comparison and a
cutoff limit prove
\(\propP\).  The general theorem then completes
Theorem~\ref{thm:exact-infiniteness}.

Section~\ref{sec:quantitative} gives an alternative final step under
exactness.  It constructs positives in two prescribed hereditary
subalgebras with arbitrarily small product, while retaining a quotient
norm floor fixed before the product error is chosen.  Two compatible
inverse perturbations preserve the two floors.  Clipping the later
comparison coefficients bounds their norms, so the small product
survives implementation and gives strong pure infiniteness directly.
Finally, for a separable nuclear algebra, the central-sequence
absorption criterion applied in Section~\ref{sec:nuclear} gives
Corollary~\ref{cor:absorption}.
